%! Semester 1, Units 1–3 — PRACTICE Final Exam (Answer Key)
% The practice form of the cumulative final. Item for item it is the parallel form of
% ../actual_test/: the same three sections, the same 20 multiple-choice targets in the same
% order, the same five short-answer computations, and the same three free-response questions
% with the same parts and point values — every context different, so working it practises the
% moves without giving away the answers. shared/tests.mk publishes its PDF to
% ../../sample_test/main.pdf, and its key to ../../sample_test_key/main.pdf.
% TIME: written for 60–120 minutes (about 2 min per MC, 2 min per short answer, 12–15 min per
% free-response question: roughly 90 minutes for a typical student; 2 hours is the cap).
% WEIGHTS: Sections I and II follow the AP exam's unit weights (Unit 1 ~45%, Unit 3 ~35%,
% Unit 2 ~20%: 11 / 9 / 5 of the 25 items); Section III is one question per unit.
% POINTS: MC 20 x 2 = 40, short answer 5 x 3 = 15, free response 3 x 15 = 45; total 100.
%
% BLUEPRINT
%   MC 1 categorical vs quantitative (digits)      marathon records                      1.1
%   MC 2 groups of unequal size                    two middle schools, walkers           1.2
%   MC 3 IQR from a stem plot                      15 daily high temperatures            1.3, 1.5
%   MC 4 shape, and mean vs median                 60 quiz scores (skewed left)          1.4, 1.5
%   MC 5 resistance                                house prices + a mansion              1.5
%   MC 6 percentile                                a newborn's birth weight              1.5
%   MC 7 every section holds a quarter             60 customers on hold                  1.6
%   MC 8 parameter vs statistic                    defective bolts                       1.7
%   MC 9 empirical rule                            sunflower heights                     1.7
%   MC10 conditional relative frequency            species x where the pet came from     2.1
%   MC11 what r says and does not say              latitude vs January temperature       2.2
%   MC12 r from R-Sq, sign from the slope          training hours vs points scored       2.4
%   MC13 influential point                         tractor age vs resale value           2.5
%   MC14 which study is an experiment              five studies                          3.1
%   MC15 name the sampling method (stratified)     four colleges                         3.2
%   MC16 name the bias and its direction           exercise, gym membership lists        3.3
%   MC17 a larger sample does not fix bias         online poll vs SRS                    3.3
%   MC18 confounding                               daily vitamins and colds              3.4
%   MC19 randomized block design                   reading program, two language groups  3.4
%   MC20 scope of inference (random sample, not assigned)   sleep and traffic tickets    3.5
%   SA21 z-score                                   5K finishing time                     1.7
%   SA22 upper outlier fence                       minutes of homework                   1.6
%   SA23 residual                                  sunflower height                      2.3
%   SA24 count the treatments (two factors)        detergent x water temperature         3.4
%   SA25 read a reshuffle dotplot                  100 reshuffles, observed 3            3.5
%   FRQ26 compare, outlier rule, resistance, empirical rule, z   earbud battery life     1.4–1.7
%   FRQ27 describe, equation + slope, residual, R-Sq, residual plot + extrapolation
%                                                  outdoor temperature vs natural gas    2.2–2.5
%   FRQ28 SRS, nonresponse + direction, units/treatments/response, random assignment, scope
%                                                  a university dining service           3.2–3.5
% ANSWERS (MC): 1C 2A 3B 4E 5D 6A 7E 8C 9D 10D 11B 12E 13C 14A 15D 16B 17A 18E 19B 20C
%
% THE DATA (every number checked)
%  MC2  East 150/500 = 0.30; West 120/300 = 0.40.
%  MC3  54 58 61 63 65 69 70 72 72 76 78 81 84 90 93: Q1 63, med 72, Q3 81, IQR 18, range 39.
%  MC4  bins of 10 points from 30: 1 1 3 6 11 20 18 (60), skewed left.
%  MC5  210 220 230 240 250: mean 230, median 230; with 890: mean 2040/6 = 340, median 235.
%  MC7  5/12/15/24/40, n = 60, so 15 customers per section.
%  MC9  N(180, 12): 168 = -1 SD -> (100% - 68%)/2 = 16%.
%  MC10 Dog 90/110 (200), Cat 126/54 (180), totals 216/164/380 -> 126/216 = 0.583.
%  MC12 n 24, xbar 30, s_x 8, ybar 62, s_y 9, r 0.70 -> b 0.7875, a 38.38, S 6.572,
%       SE(b) 0.1713, T 4.60; SE(a) 5.311, T 7.23; R-Sq 49.0%, adj 46.7%.
%  MC13 x 2..12, y 48 45 46 41 40 36 37 32 31 27 26: b -2.25, r -0.987. With (24, 44): b -0.38,
%       r -0.295.
%  SA21 (23 - 28)/4 = -1.25.   SA22 12/30/41/52/98: IQR 22, 52 + 33 = 85.
%  SA23 42 + 5.5(6) = 75; 79 - 75 = 4.   SA24 2 x 4 = 8.
%  SA25 100 reshuffles, step 1: -5:1 -4:2 -3:6 -2:12 -1:18 0:22 1:18 2:12 3:6 4:2 5:1
%       (sum 100). At or above 3: 6 + 2 + 1 = 9.
%  FRQ26 Brand A 4.0/5.2/6.0/6.8/8.1 (IQR 1.6, fences 2.8 and 9.2: no outliers), about
%       N(6.0, 1.2). Brand B 5.5/7.0/7.6/8.2, whisker to 9.8, outlier 11.0 (IQR 1.2, fence 10.0).
%       (d) 3.6 = 6.0 - 2(1.2) -> 2.5%. (e) (8.4-6.0)/1.2 = 2.0; (8.8-7.6)/0.8 = 1.5.
%  FRQ27 16 days, outdoor temperature x (F): 12 15 18 21 24 27 30 33 36 39 42 45 48 51 54 57
%        gas used y (therms): 9.9 9.0 9.2 8.2 8.1 7.0 6.9 5.9 6.5 5.3 5.3 4.6 3.7 3.7 3.0 2.6
%        yhat = 11.69 - 0.1597x, r -0.992, R-Sq 98.3%, adj 98.2%, S 0.3075, SE(b) 0.00556,
%        T -28.72, SE(a) 0.2066, T 56.57. Residuals: .13 -.29 .38 -.14 .24 -.38 .00 -.52 .56
%        -.16 .32 .10 -.33 .15 -.07 .01 (random scatter). (36, 6.5): yhat 5.94, residual 0.56
%        with the printed coefficients. At 90 F: -2.68 therms.
% NAMESTRIP: no name line — mirrors the blank (a graded page; the name is on the cover sheet only).
% NO teachernote here — scoring guidance belongs in the rubrics, never in a key.
% GENERATED FROM the blank.
% PAGINATION: each item carries a \boxguard sized to its height so no stem, figure or option
% list splits across a page; each section strip is guarded to travel with its first item.
% NO SKETCH ITEMS: every display is pre-drawn in TikZ and read.
% PAGE LOCKSTEP with ../../tests/practice_test/: every \writelines{n} here is answered by
% exactly n \ansline{} there, each terse enough not to wrap; every work block is byte-identical.
% Nothing else differs except the preamble, the header, and the \ans{} marks on the correct MC
% options.
\documentclass[12pt]{article}
\usepackage{apstats-article}
\usepackage{apstats-key}   % pulls in -boxes; do NOT also load -boxes
\usepackage{grading}
\gradingassignment{apstats-s01-practice-final}
\gradingcourse{advanced_statistics}

\begin{document}

\pageheader{Semester 1 \quad Units 1--3}{Practice Final Exam --- Answer Key}
{\small\itshape Write each final answer in its box. Only what is inside a box is graded.}\par

\hfill\small Score:\;\blank{1.3cm}\;/\;100\normalsize
\vspace{0.06in}

\boxguard[8]
\begin{headlinebox}{sky}
\small
\textbf{The final allows up to 2 hours.} \textbf{Section I:} 20 multiple-choice questions, 40
points. \textbf{Section II:} 5 short-answer questions, 15 points. \textbf{Section III:} 3
free-response questions, 45 points. The exam covers Units 1, 2, and 3.
\end{headlinebox}

\vspace{0.08in}

\boxguard[10]
\begin{headlinebox}{goldbg}
\small
\textbf{Section I --- Multiple Choice} (2 points each). Write the letter of the single best answer in its Letter box. Read all
five options --- several of them name a real idea from the course, but only one fits \emph{this}
question.
\end{headlinebox}

\vspace{0.10in}

\begin{enumerate}[leftmargin=1.9em, itemsep=0.13in]

  % ── MC 1 ── 1.1
  \boxguard[8]
  \item A marathon records five variables for each runner: \emph{bib number}, \emph{age}
        (years), \emph{finishing time} (minutes), \emph{shirt size}, and \emph{number of
        marathons run before}.

        How many of these five variables are \textbf{quantitative}?\hfill\gradedright{\textit{Letter:}~\gradedblank[id=t1, type=choice, points=2, answer=C]{0.45in}{0.2in}}\rule[-0.2in]{0pt}{0.48in}
        \par\vspace{3pt}\noindent
        \hspace{1.2em}\textbf{(A)} 1 \hfill \textbf{(B)} 2 \hfill
        \textbf{(C)} \ans{3} \hfill \textbf{(D)} 4 \hfill \textbf{(E)} 5

  % ── MC 2 ── 1.2
  \boxguard[11]
  \item East Middle School has 500 students, and 150 of them walk to school. West Middle
        School has 300 students, and 120 of them walk to school. Which statement is best
        supported?
        \begin{enumerate}[label=\textbf{(\Alph*)}, leftmargin=2.2em, itemsep=2pt]
          \item \ans{Walking is more common at West: 40\% of its students walk, compared with}\\
                \ans{30\% of East's.}
          \item Walking is more common at East, because 150 students is more than 120 students.
          \item Walking is equally common, because each school has more than 100 walkers.
          \item No comparison is possible, because the schools are different sizes.
          \item Most West students walk to school.\hfill\gradedright{\textit{Letter:}~\gradedblank[id=t2, type=choice, points=2, answer=A]{0.45in}{0.2in}}\rule[-0.2in]{0pt}{0.48in}
        \end{enumerate}

  % ── MC 3 ── 1.3, 1.5
  \boxguard[13]
  \item The stem plot shows the high temperature (degrees Fahrenheit) in a town on each of
        15 days in May.

        \vspace{3pt}
        {\centering\small
        \renewcommand{\arraystretch}{1.05}
        \begin{tabular}{@{} r | l @{}}
        5 & 4 \ 8 \\
        6 & 1 \ 3 \ 5 \ 9 \\
        7 & 0 \ 2 \ 2 \ 6 \ 8 \\
        8 & 1 \ 4 \\
        9 & 0 \ 3 \\
        \end{tabular}\qquad
        \raisebox{0.28in}{\footnotesize Key: $5 \mid 4$ means 54$^\circ$F}\par}
        \vspace{2pt}

        What is the \textbf{interquartile range} of the temperatures?\hfill\gradedright{\textit{Letter:}~\gradedblank[id=t3, type=choice, points=2, answer=B]{0.45in}{0.2in}}\rule[-0.2in]{0pt}{0.48in}
        \par\vspace{3pt}\noindent
        \hspace{1.2em}\textbf{(A)} 9 \hfill \textbf{(B)} \ans{18} \hfill
        \textbf{(C)} 21 \hfill \textbf{(D)} 39 \hfill \textbf{(E)} 72

  % ── MC 4 ── 1.4, 1.5
  \boxguard[18]
  \item The histogram shows the scores of 60 students on a 100-point quiz.

        \vspace{3pt}
        {\centering
        \begin{tikzpicture}[x=0.135cm, y=0.125cm]
          \foreach \y in {0,5,10,15,20} {
            \draw[linegray, very thin] (30,\y) -- (102,\y);
            \node[left, font=\tiny] at (30,\y) {\y};}
          \foreach \l/\c in {30/1, 40/1, 50/3, 60/6, 70/11, 80/20, 90/18}
            \draw[fill=skymid] (\l,0) rectangle (\l+10,\c);
          \draw[thick] (30,0) -- (102,0);
          \draw[thick] (30,0) -- (30,22);
          \foreach \x in {30,40,...,100} { \node[below, font=\tiny] at (\x,0) {\x}; }
          \node[font=\tiny] at (66,-5.0) {Quiz score (points)};
          \node[rotate=90, font=\tiny] at (23.5,11) {Students};
        \end{tikzpicture}\par}
        \vspace{1pt}

        Which statement is most likely true?\hfill\gradedright{\textit{Letter:}~\gradedblank[id=t4, type=choice, points=2, answer=E]{0.45in}{0.2in}}\rule[-0.2in]{0pt}{0.48in}
        \begin{enumerate}[label=\textbf{(\Alph*)}, leftmargin=2.2em, itemsep=2pt]
          \item The distribution is skewed right, so the mean is likely greater than the median.
          \item The distribution is skewed right, so the mean is likely less than the median.
          \item The distribution is roughly symmetric, so the mean and median are about equal.
          \item The distribution is skewed left, so the mean is likely greater than the median.
          \item \ans{The distribution is skewed left, so the mean is likely less than the median.}
        \end{enumerate}

  % ── MC 5 ── 1.5
  \boxguard[11]
  \item Five houses on one street sold for 210, 220, 230, 240, and 250 thousand dollars. A
        sixth house, a mansion, then sold for 890 thousand dollars and is added to the data
        set. What happens to the mean and the median?
        \begin{enumerate}[label=\textbf{(\Alph*)}, leftmargin=2.2em, itemsep=2pt]
          \item Both increase by the same amount.
          \item The mean increases and the median decreases.
          \item The median increases from 230 to 340; the mean only from 230 to 235.
          \item \ans{The mean increases from 230 to 340; the median only from 230 to 235.}
          \item Neither one changes.\hfill\gradedright{\textit{Letter:}~\gradedblank[id=t5, type=choice, points=2, answer=D]{0.45in}{0.2in}}\rule[-0.2in]{0pt}{0.48in}
        \end{enumerate}

  % ── MC 6 ── 1.5
  \boxguard[10]
  \item A newborn's birth weight is at the 30th percentile for babies born at one hospital.
        Which statement is correct?
        \begin{enumerate}[label=\textbf{(\Alph*)}, leftmargin=2.2em, itemsep=2pt]
          \item \ans{30\% of the babies born at the hospital weighed as much as this baby or less.}
          \item This baby weighs 30\% of the mean birth weight at the hospital.
          \item 30\% of the babies born at the hospital weighed more than this baby.
          \item This baby weighed 30 ounces.
          \item This baby is 30\% heavier than the lightest baby.\hfill\gradedright{\textit{Letter:}~\gradedblank[id=t6, type=choice, points=2, answer=A]{0.45in}{0.2in}}\rule[-0.2in]{0pt}{0.48in}
        \end{enumerate}

  % ── MC 7 ── 1.6
  \boxguard[16]
  \item The boxplot shows the number of minutes each of 60 customers waited on hold for a
        help line.

        \vspace{3pt}
        {\centering
        \begin{tikzpicture}[x=0.27cm, y=1cm]
          \draw[thick] (0,0) -- (44,0);
          \foreach \x in {0,10,...,40} {
            \draw (\x,-0.07) -- (\x,0.07); \node[below, font=\tiny] at (\x,-0.05) {\x};}
          \draw (5,0.45) -- (12,0.45);
          \draw (24,0.45) -- (40,0.45);
          \draw[fill=skymid] (12,0.25) rectangle (24,0.65);
          \draw[thick] (15,0.25) -- (15,0.65);
          \draw (5,0.33) -- (5,0.57);
          \draw (40,0.33) -- (40,0.57);
          \foreach \x in {5,12,15,24,40} \node[above, font=\tiny] at (\x,0.66) {\x};
          \node[font=\tiny] at (22,-0.62) {Minutes on hold};
        \end{tikzpicture}\par}
        \vspace{1pt}

        Which statement is true?\hfill\gradedright{\textit{Letter:}~\gradedblank[id=t7, type=choice, points=2, answer=E]{0.45in}{0.2in}}\rule[-0.2in]{0pt}{0.48in}
        \begin{enumerate}[label=\textbf{(\Alph*)}, leftmargin=2.2em, itemsep=2pt]
          \item More customers waited 24 to 40 minutes than 12 to 15, because that whisker is
                longer.
          \item About half of the customers waited more than 24 minutes.
          \item The mean wait was 15 minutes.
          \item No customer waited more than 24 minutes.
          \item \ans{About 15 customers waited 12 to 15 minutes, and about 15 waited 24 to 40.}
        \end{enumerate}

  % ── MC 8 ── 1.7
  \boxguard[11]
  \item A factory makes 80{,}000 bolts a day. An inspector tests a random sample of 250 of
        today's bolts and finds that 2.4\% of them are defective. Which statement is correct?
        \begin{enumerate}[label=\textbf{(\Alph*)}, leftmargin=2.2em, itemsep=2pt]
          \item 2.4\% is a parameter, because it is a percent.
          \item 2.4\% is a parameter, because the sample was chosen at random.
          \item \ans{2.4\% is a statistic; the percent of all 80{,}000 of today's bolts that are}\\
                \ans{defective is the parameter.}
          \item 250 is the statistic and 80{,}000 is the parameter.
          \item 2.4\% is both, because the sample is random.\hfill\gradedright{\textit{Letter:}~\gradedblank[id=t8, type=choice, points=2, answer=C]{0.45in}{0.2in}}\rule[-0.2in]{0pt}{0.48in}
        \end{enumerate}

  % ── MC 9 ── 1.7
  \boxguard[8]
  \item The heights of the sunflowers in a field are approximately normal, with mean 180
        centimeters and standard deviation 12 centimeters.

        About what percent of the sunflowers are shorter than 168 centimeters?\hfill\gradedright{\textit{Letter:}~\gradedblank[id=t9, type=choice, points=2, answer=D]{0.45in}{0.2in}}\rule[-0.2in]{0pt}{0.48in}
        \par\vspace{3pt}\noindent
        \hspace{1.2em}\textbf{(A)} 2.5\% \hfill \textbf{(B)} 5\% \hfill
        \textbf{(C)} 13.5\% \hfill \textbf{(D)} \ans{16\%} \hfill \textbf{(E)} 32\%

  % ── MC 10 ── 2.1
  \boxguard[13]
  \item A veterinarian classified 380 pets by species and by where the owner got the pet.

        \vspace{2pt}
        {\centering\small
        \renewcommand{\arraystretch}{1.15}
        \begin{tabular}{@{} l c c c @{}}
        \rowcolor{sky}
         & \textbf{From a shelter} & \textbf{From a breeder} & \textbf{Total} \\
        \midrule
        Dog   & 90  & 110 & 200 \\
        Cat   & 126 & 54  & 180 \\
        \midrule
        Total & 216 & 164 & 380 \\
        \end{tabular}\par}
        \vspace{4pt}

        What proportion of the pets that came \textbf{from a shelter} are cats?\hfill\gradedright{\textit{Letter:}~\gradedblank[id=t10, type=choice, points=2, answer=D]{0.45in}{0.2in}}\rule[-0.2in]{0pt}{0.48in}
        \par\vspace{3pt}\noindent
        \hspace{1.2em}\textbf{(A)} 0.332 \hfill \textbf{(B)} 0.474 \hfill
        \textbf{(C)} 0.568 \hfill \textbf{(D)} \ans{0.583} \hfill \textbf{(E)} 0.70

  % ── MC 11 ── 2.2
  \boxguard[11]
  \item For 35 cities, the correlation between latitude (degrees north) and mean January
        temperature (degrees Fahrenheit) is $r = -0.71$. Which statement is correct?
        \begin{enumerate}[label=\textbf{(\Alph*)}, leftmargin=2.2em, itemsep=2pt]
          \item 71\% of the cities that are farther north are colder in January.
          \item \ans{There is a fairly strong, negative linear association: cities farther north}\\
                \ans{tend to have lower January temperatures.}
          \item Latitude is the only cause of a city's January temperature.
          \item If temperature were measured in degrees Celsius, $r$ would change.
          \item Temperature falls 0.71$^\circ$F for each degree north.\hfill\gradedright{\textit{Letter:}~\gradedblank[id=t11, type=choice, points=2, answer=B]{0.45in}{0.2in}}\rule[-0.2in]{0pt}{0.48in}
        \end{enumerate}

  % ── MC 12 ── 2.4
  \boxguard[11]
  \item A swim coach regressed the points each of 24 swimmers scored this season on the hours
        the swimmer trained per month.

        \vspace{2pt}
        {\centering\footnotesize\ttfamily
        \begin{tabular}{@{} l r r r r @{}}
        Predictor & Coef & SE Coef & T & P \\ \hline
        Constant & 38.38 & 5.311 & 7.23 & 0.000 \\
        Hours & 0.7875 & 0.1713 & 4.60 & 0.000 \\
        \end{tabular}\par
        S = 6.572 \quad R-Sq = 49.0\% \quad R-Sq(adj) = 46.7\%\par}
        \vspace{2pt}

        What is the correlation between training hours and points scored?\hfill\gradedright{\textit{Letter:}~\gradedblank[id=t12, type=choice, points=2, answer=E]{0.45in}{0.2in}}\rule[-0.2in]{0pt}{0.48in}
        \par\vspace{3pt}\noindent
        \hspace{1.2em}\textbf{(A)} $-0.70$ \hfill \textbf{(B)} $-0.49$ \hfill
        \textbf{(C)} 0.24 \hfill \textbf{(D)} 0.49 \hfill \textbf{(E)} \ans{0.70}

  % ── MC 13 ── 2.5
  \boxguard[20]
  \item The scatterplot shows the age (years) and the resale value (thousands of dollars) of
        12 tractors. A least-squares line is fit to all 12.

        \vspace{3pt}
        {\centering
        \begin{tikzpicture}[xscale=0.30, yscale=0.105]
          \foreach \y in {30,40,50} {\draw[linegray, very thin] (0,\y) -- (26,\y);
            \node[left, font=\tiny] at (0,\y) {\y};}
          \draw[->, thick] (0,22) -- (27,22);
          \draw[->, thick] (0,22) -- (0,53);
          \foreach \x in {0,5,10,15,20,25} {\draw (\x,21.4) -- (\x,22.6);
            \node[below, font=\tiny] at (\x,21.6) {\x};}
          \foreach \x/\y in {2/48, 3/45, 4/46, 5/41, 6/40, 7/36, 8/37, 9/32, 10/31, 11/27,
                             12/26, 24/44}
            {\fill[navy] (\x,\y) circle[x radius=0.27, y radius=0.76];}
          \node[font=\tiny] at (13.5,15.2) {Age (years)};
          \node[rotate=90, font=\tiny] at (-3.4,37) {Value (\$1000s)};
        \end{tikzpicture}\par}
        \vspace{1pt}

        If the 24-year-old tractor were removed, what would happen?\hfill\gradedright{\textit{Letter:}~\gradedblank[id=t13, type=choice, points=2, answer=C]{0.45in}{0.2in}}\rule[-0.2in]{0pt}{0.48in}
        \begin{enumerate}[label=\textbf{(\Alph*)}, leftmargin=2.2em, itemsep=2pt]
          \item Nothing: one point out of 12 cannot change a least-squares line.
          \item The slope would become less steep, and $r$ would move closer to 0.
          \item \ans{The slope would become steeper, and $r$ would move closer to $-1$.}
          \item The slope would stay the same, and $r$ would become positive.
          \item The slope would become steeper, and $r$ would move closer to 0.
        \end{enumerate}

  % ── MC 14 ── 3.1
  \boxguard[14]
  \item Which of these studies is an \textbf{experiment}?
        \begin{enumerate}[label=\textbf{(\Alph*)}, leftmargin=2.2em, itemsep=2pt]
          \item \ans{A gardener randomly chooses 20 of her 40 tomato plants to get a new fertilizer,
                gives the other 20 her usual one, and compares the yields.}
          \item A dentist uses patient records to compare the cavities of patients who floss
                daily with those of patients who do not.
          \item A researcher surveys 500 adults and compares the stress levels of pet owners
                with those of adults who have no pet.
          \item A fitness tracker logs 8{,}000 users and compares early risers with night owls.
          \item A school asks 300 randomly chosen parents about bus routes.\hfill\gradedright{\textit{Letter:}~\gradedblank[id=t14, type=choice, points=2, answer=A]{0.45in}{0.2in}}\rule[-0.2in]{0pt}{0.48in}
        \end{enumerate}

  % ── MC 15 ── 3.2
  \boxguard[13]
  \item A university has four colleges: arts, business, engineering, and science. To study
        student debt, it takes a simple random sample of 50 students from each of the four
        colleges. What type of sample is this?
        \begin{enumerate}[label=\textbf{(\Alph*)}, leftmargin=2.2em, itemsep=2pt]
          \item A simple random sample, because the students were chosen at random.
          \item A cluster sample, because the students are grouped into colleges.
          \item A convenience sample, because the students were already at the university.
          \item \ans{A stratified random sample, because a random sample was taken from every}\\
                \ans{college.}
          \item A voluntary response sample, because students can decline.\hfill\gradedright{\textit{Letter:}~\gradedblank[id=t15, type=choice, points=2, answer=D]{0.45in}{0.2in}}\rule[-0.2in]{0pt}{0.48in}
        \end{enumerate}

  % ── MC 16 ── 3.3
  \boxguard[13]
  \item To estimate the proportion of adults in a town who exercise at least three times a
        week, a researcher takes a random sample of names from the membership lists of the
        town's gyms. Which is the main problem with this plan?
        \begin{enumerate}[label=\textbf{(\Alph*)}, leftmargin=2.2em, itemsep=2pt]
          \item Voluntary response bias, so the estimate is likely too low.
          \item \ans{Undercoverage: adults who do not belong to a gym cannot be chosen,}\\
                \ans{so the estimate is likely too high.}
          \item Response bias, because the wording of the question is leading.
          \item There is no problem, because the sample was chosen at random.
          \item The sample is too small, so the estimate is likely too low.\hfill\gradedright{\textit{Letter:}~\gradedblank[id=t16, type=choice, points=2, answer=B]{0.45in}{0.2in}}\rule[-0.2in]{0pt}{0.48in}
        \end{enumerate}

  % ── MC 17 ── 3.3
  \boxguard[14]
  \item An online magazine posts a poll asking whether readers support a new school dress
        code; 40{,}000 readers respond, and 81\% of them oppose it. A research firm asks a
        simple random sample of 800 adults in the country; 52\% oppose it. Which result is the
        more trustworthy estimate for the country's adults?
        \begin{enumerate}[label=\textbf{(\Alph*)}, leftmargin=2.2em, itemsep=2pt]
          \item \ans{The SRS of 800: a larger sample does not remove the bias of a voluntary}\\
                \ans{response sample.}
          \item The online poll, because 40{,}000 is a much larger sample than 800.
          \item The online poll, because its readers chose to answer.
          \item Neither result, because neither one is a census.
          \item Both equally, because both asked the same question.\hfill\gradedright{\textit{Letter:}~\gradedblank[id=t17, type=choice, points=2, answer=A]{0.45in}{0.2in}}\rule[-0.2in]{0pt}{0.48in}
        \end{enumerate}

  % ── MC 18 ── 3.4
  \boxguard[14]
  \item A health magazine reports that adults who choose to take a daily vitamin get fewer
        colds, on average, than adults who do not. The magazine concludes that the vitamin
        prevents colds.

        Which is the best criticism?\hfill\gradedright{\textit{Letter:}~\gradedblank[id=t18, type=choice, points=2, answer=E]{0.45in}{0.2in}}\rule[-0.2in]{0pt}{0.48in}
        \begin{enumerate}[label=\textbf{(\Alph*)}, leftmargin=2.2em, itemsep=2pt]
          \item The study had no placebo, so the numbers of colds cannot be compared.
          \item The adults should have been chosen by a stratified random sample.
          \item A count of colds is categorical, so the two groups cannot be compared.
          \item None is needed: the adults who took the vitamin had fewer colds.
          \item \ans{Adults chose whether to take the vitamin, so health habits are confounded with}\\
                \ans{it: health-conscious adults may take vitamins and get fewer colds anyway.}
        \end{enumerate}

  % ── MC 19 ── 3.4
  \boxguard[13]
  \item A district tests a new reading program against its current one with 80 third-graders:
        40 are native English speakers, and 40 are English learners. Language background is
        known to affect reading scores. Which plan is a \textbf{randomized block design}?
        \begin{enumerate}[label=\textbf{(\Alph*)}, leftmargin=2.2em, itemsep=2pt]
          \item Give the new program to every English learner and the current one to every
                native speaker.
          \item \ans{Within each language group separately, randomly assign 20 students to}\\
                \ans{the new program and 20 to the current one.}
          \item Randomly assign 40 of the 80 students to the new program and the rest to the
                current one.
          \item Let each student's parents choose the program.
          \item Randomly select 40 students and drop the rest.\hfill\gradedright{\textit{Letter:}~\gradedblank[id=t19, type=choice, points=2, answer=B]{0.45in}{0.2in}}\rule[-0.2in]{0pt}{0.48in}
        \end{enumerate}

  % ── MC 20 ── 3.5
  \boxguard[13]
  \item A researcher takes a random sample of 600 adults from a state's list of licensed
        drivers and asks each one how many hours they sleep per night and how many traffic
        tickets they have received. Drivers who sleep less have significantly more tickets.
        Which conclusion is justified?
        \begin{enumerate}[label=\textbf{(\Alph*)}, leftmargin=2.2em, itemsep=2pt]
          \item Sleeping less causes the state's drivers to get more tickets.
          \item Sleeping less causes these 600 drivers, but no others, to get more tickets.
          \item \ans{Among the state's licensed drivers, those who sleep less tend to have more}\\
                \ans{tickets, but no cause can be claimed.}
          \item Nothing can be concluded, because the drivers were not volunteers.
          \item Nothing, because 600 drivers is too few.\hfill\gradedright{\textit{Letter:}~\gradedblank[id=t20, type=choice, points=2, answer=C]{0.45in}{0.2in}}\rule[-0.2in]{0pt}{0.48in}
        \end{enumerate}

\end{enumerate}

\vspace{0.14in}

\boxguard[22]
\begin{headlinebox}{goldbg}
\small
\textbf{Section II --- Short Answer} (3 points each). Work each problem in the space under it.
Write only the final number in its box. Only what is inside a box is graded.
\end{headlinebox}

\vspace{0.10in}

\begin{enumerate}[leftmargin=1.9em, itemsep=0.16in, resume]

  % ── SA 21 ── 1.7
  \boxguard[11]
  \item The finishing times in a 5-kilometer race are approximately normal, with mean 28
        minutes and standard deviation 4 minutes. One runner finished in 23 minutes. Find the
        $z$-score of this runner's time.\gradedbelow[r]{0.62in}{$z = $~\gradedblank[id=t21, type=short, points=3, answer=-1.25, check=numeric]{1.1in}{0.34in}}
        \begin{work}
          z &= \frac{23 - 28}{4} = \frac{-5}{4} = -1.25
        \end{work}
        \vspace{0.62in}

  % ── SA 22 ── 1.6
  \boxguard[11]
  \item The five-number summary of the minutes 40 students spent on homework last night is
        12, 30, 41, 52, 98. By the $1.5 \times \text{IQR}$ rule, a value \emph{above} what
        number is an outlier?\gradedbelow[r]{0.62in}{Above~\gradedblank[id=t22, type=short, points=3, answer=85, check=numeric]{1.1in}{0.34in}}
        \begin{work}
          \text{IQR} &= 52 - 30 = 22 \\
          Q_3 + 1.5(\text{IQR}) &= 52 + 1.5(22) = 85
        \end{work}
        \vspace{0.50in}

  % ── SA 23 ── 2.3
  \boxguard[11]
  \item The line $\hat{y} = 42 + 5.5x$ predicts a sunflower's height $\hat{y}$ (centimeters)
        from the number of weeks $x$ since it sprouted. A sunflower that sprouted 6 weeks ago
        is 79 centimeters tall. Find the residual for this sunflower.\gradedbelow[r]{0.62in}{Residual $=$~\gradedblank[id=t23, type=short, points=3, answer=4, check=numeric]{1.1in}{0.34in}}
        \begin{work}
          \hat{y} &= 42 + 5.5(6) = 75 \\
          \text{residual} &= y - \hat{y} = 79 - 75 = 4
        \end{work}
        \vspace{0.50in}

  % ── SA 24 ── 3.4
  \boxguard[10]
  \item An experiment on stain removal uses two factors: the brand of detergent (A or B) and
        the water temperature (cold, cool, warm, or hot). Every combination of detergent and
        temperature is used. How many treatments are there?\gradedbelow[r]{0.50in}{Treatments:~\gradedblank[id=t24, type=short, points=3, answer=8, check=numeric]{1.1in}{0.34in}}
        \begin{work}
          2 \text{ brands of detergent} \times 4 \text{ water temperatures} &= 8 \text{ treatments}
        \end{work}
        \vspace{0.60in}

  % ── SA 25 ── 3.5
  \boxguard[20]
  \item In an experiment, the mean score of the treatment group was \textbf{3 points higher}
        than the mean score of the control group. A computer reshuffled the scores into two new
        groups 100 times, as if the treatment made no difference. The dotplot shows the 100
        simulated differences. How many of the 100 reshuffles gave a difference of 3 points or
        more?
        \begin{center}
        \begin{tikzpicture}[x=0.95cm, y=0.105cm]
          \draw[thick] (-5.6,0) -- (5.6,0);
          \foreach \x in {-5,...,5} {
            \draw (\x,0) -- (\x,-1.4); \node[below, font=\tiny] at (\x,-1.2) {$\x$};}
          \foreach \x/\n in {-5/1,-4/2,-3/6,-2/12,-1/18,0/22,1/18,2/12,3/6,4/2,5/1}
            \foreach \k in {1,...,\n} \fill[navy] (\x,\k) circle (1.3pt);
          \node[font=\tiny] at (0,-7.2) {Simulated difference in mean score (points), treatment $-$ control};
          \node[font=\scriptsize\bfseries] at (0,27) {100 random reshuffles};
        \end{tikzpicture}
        \end{center}
        \vspace{-6pt}
        \begin{work}
          6 + 2 + 1 &= 9
        \end{work}
        \hfill Reshuffles at 3 or more:~\gradedblank[id=t25, type=short, points=3, answer=9, check=numeric]{1.1in}{0.34in}

\end{enumerate}

\vspace{0.14in}

\boxguard[34]
\begin{headlinebox}{goldbg}
\small
\textbf{Section III --- Free Response} (15 points each). Each part has its own box: write your
final answer to the part inside it. Show calculations in the space above the box. Answer
\emph{in context} --- name the people or things studied, the variables, and their units.
\end{headlinebox}

\vspace{0.10in}

\begin{enumerate}[leftmargin=1.9em, itemsep=0.16in, resume]

  % ── FRQ 26 ── Unit 1
  \boxguard[32]
  \item \textbf{(15 points)} A reviewer measured the battery life (hours) of 40 pairs of
        Brand A earbuds and 40 pairs of Brand B earbuds.

        \vspace{4pt}
        \noindent
        \begin{minipage}[c]{0.55\linewidth}\centering
        \begin{tikzpicture}[x=0.92cm, y=1.15cm]
          \draw[thick] (3,0) -- (12,0);
          \foreach \x in {3,4,...,12} {
            \draw (\x,-0.07) -- (\x,0.07); \node[below, font=\tiny] at (\x,-0.05) {\x};}
          % Brand A 4.0/5.2/6.0/6.8/8.1
          \draw (4.0,1.15) -- (5.2,1.15); \draw (6.8,1.15) -- (8.1,1.15);
          \draw[fill=skymid] (5.2,0.95) rectangle (6.8,1.35); \draw[thick] (6.0,0.95) -- (6.0,1.35);
          \draw (4.0,1.03) -- (4.0,1.27); \draw (8.1,1.03) -- (8.1,1.27);
          \node[font=\tiny, anchor=south west] at (3,1.38) {Brand A};
          % Brand B 5.5/7.0/7.6/8.2, whisker 9.8, outlier 11.0
          \draw (5.5,0.45) -- (7.0,0.45); \draw (8.2,0.45) -- (9.8,0.45);
          \draw[fill=goldbg] (7.0,0.25) rectangle (8.2,0.65); \draw[thick] (7.6,0.25) -- (7.6,0.65);
          \draw (5.5,0.33) -- (5.5,0.57); \draw (9.8,0.33) -- (9.8,0.57);
          \fill[navy] (11.0,0.45) circle (1.6pt);
          \node[font=\tiny, anchor=south west] at (3,0.68) {Brand B};
          \node[font=\tiny] at (7.5,-0.55) {Battery life (hours)};
        \end{tikzpicture}
        \end{minipage}\hfill
        \begin{minipage}[c]{0.42\linewidth}\centering\footnotesize
        \renewcommand{\arraystretch}{1.15}
        \setlength{\tabcolsep}{3.5pt}
        \begin{tabular}{@{} l c c c c c @{}}
        \rowcolor{sky}
         & \textbf{Min} & $\boldsymbol{Q_1}$ & \textbf{Med} & $\boldsymbol{Q_3}$ & \textbf{Max} \\
        \midrule
        Brand A & 4.0 & 5.2 & 6.0 & 6.8 & 8.1 \\
        Brand B & 5.5 & 7.0 & 7.6 & 8.2 & 11.0 \\
        \end{tabular}
        \end{minipage}

        \vspace{6pt}
        \textbf{(a)} Compare the distributions of battery life for Brand A and Brand B
        earbuds.\\[4pt]
        \gradedhere{\gradedblank[id=t26a, type=rubric, points=3, rubric=../rubrics/apstats-s01-practice-final-t26a.yaml]{\linewidth+6pt}{\gradedlinesheight{3}}}\ansline{Brand B batteries typically last longer: median 7.6 h vs.\ 6.0 h for Brand A.}\\
        \ansline{Brand A battery lives vary more: IQR 1.6 h vs.\ 1.2 h for Brand B.}\\
        \ansline{Brand B has one high outlier (11.0 h); Brand A has no outliers.}\\
        \ansline{}\\
        \ansline{}\\
        \ansline{}\\
        \ansline{}\\
        \ansline{}

        \vspace{6pt}
        \boxguard[20]
        \textbf{(b)} Use the $1.5 \times \text{IQR}$ rule to show that the Brand B battery life
        of 11.0 hours is an outlier.
        \begin{work}
          \text{IQR} &= 8.2 - 7.0 = 1.2 \text{ hours} \\
          Q_3 + 1.5(\text{IQR}) &= 8.2 + 1.5(1.2) = 10.0 \text{ hours}
        \end{work}
        \strut\gradedhere{\gradedblank[id=t26b, type=rubric, points=3, rubric=../rubrics/apstats-s01-practice-final-t26b.yaml]{\linewidth+6pt}{\gradedlinesheight{3}}}\ansline{11.0 hours is greater than the upper fence of 10.0 hours, so the Brand B battery}\\
        \ansline{life of 11.0 hours is an outlier.}\\
        \ansline{}\\
        \ansline{}\\
        \ansline{}\\
        \ansline{}\\
        \ansline{}\\
        \ansline{}

        \vspace{6pt}
        \boxguard[16]
        \textbf{(c)} The reviewer finds that the 11.0 was a recording error: that pair of
        Brand B earbuds really lasted 8.0 hours. Without calculating, explain how correcting it
        changes the mean and the median of the Brand B battery lives.\\[4pt]
        \gradedhere{\gradedblank[id=t26c, type=rubric, points=3, rubric=../rubrics/apstats-s01-practice-final-t26c.yaml]{\linewidth+6pt}{\gradedlinesheight{3}}}\ansline{The mean decreases (by $3/40 = 0.075$ h): it uses every value, so it is not resistant.}\\
        \ansline{The median stays 7.6 h: 8.0 h is still above the median, so the middle of the}\\
        \ansline{ordered Brand B battery lives does not move. The median is resistant.}\\
        \ansline{}\\
        \ansline{}\\
        \ansline{}\\
        \ansline{}\\
        \ansline{}

        \vspace{6pt}
        \boxguard[16]
        \textbf{(d)} Brand A battery lives are approximately normal, with mean 6.0 hours and
        standard deviation 1.2 hours. About what percent of Brand A earbuds last \emph{less
        than} 3.6 hours? Explain.\\[4pt]
        \gradedhere{\gradedblank[id=t26d, type=rubric, points=3, rubric=../rubrics/apstats-s01-practice-final-t26d.yaml]{\linewidth+6pt}{\gradedlinesheight{3}}}\ansline{3.6 h $= 6.0 - 2(1.2)$, two standard deviations below the mean. About 95\% of}\\
        \ansline{Brand A battery lives are within 2 SDs of the mean, leaving 5\% outside, half on}\\
        \ansline{each side: about 2.5\% of Brand A earbuds last less than 3.6 hours.}\\
        \ansline{}\\
        \ansline{}\\
        \ansline{}\\
        \ansline{}\\
        \ansline{}

        \vspace{6pt}
        \boxguard[22]
        \textbf{(e)} After the correction, Brand B battery lives have mean 7.6 hours and
        standard deviation 0.8 hours. Which is more unusual \emph{for its own brand}: a
        Brand A battery life of 8.4 hours or a Brand B battery life of 8.8 hours? Justify with
        $z$-scores.
        \begin{work}
          z_{\text{Brand A}} &= \frac{8.4 - 6.0}{1.2} = 2.0 \\[3pt]
          z_{\text{Brand B}} &= \frac{8.8 - 7.6}{0.8} = 1.5
        \end{work}
        \strut\gradedhere{\gradedblank[id=t26e, type=rubric, points=3, rubric=../rubrics/apstats-s01-practice-final-t26e.yaml]{\linewidth+6pt}{\gradedlinesheight{3}}}\ansline{The Brand A battery life of 8.4 h is more unusual: it is 2.0 standard deviations}\\
        \ansline{above the Brand A mean, but 8.8 h is only 1.5 SDs above the Brand B mean.}\\
        \ansline{}\\
        \ansline{}\\
        \ansline{}\\
        \ansline{}\\
        \ansline{}\\
        \ansline{}

  % ── FRQ 27 ── Unit 2
  \boxguard[36]
  \item \textbf{(15 points)} A homeowner recorded the mean outdoor temperature (degrees
        Fahrenheit) and the natural gas used to heat the house (therms) on 16 winter days. The
        scatterplot, computer output, and residual plot are below.

        \vspace{4pt}
        \noindent
        \begin{minipage}[b]{0.49\linewidth}\centering
        \begin{tikzpicture}[xscale=0.125, yscale=0.40]
          \foreach \y in {2,4,6,8,10} {\draw[linegray, very thin] (8,\y) -- (60,\y);
            \node[left, font=\tiny] at (8,\y) {\y};}
          \draw[->, thick] (8,1) -- (62,1);
          \draw[->, thick] (8,1) -- (8,10.8);
          \foreach \x in {10,20,30,40,50,60} {\draw (\x,0.87) -- (\x,1.13);
            \node[below, font=\tiny] at (\x,0.95) {\x};}
          \foreach \x/\y in {12/9.9, 15/9.0, 18/9.2, 21/8.2, 24/8.1, 27/7.0, 30/6.9, 33/5.9,
                             36/6.5, 39/5.3, 42/5.3, 45/4.6, 48/3.7, 51/3.7, 54/3.0, 57/2.6}
            {\fill[navy] (\x,\y) circle[x radius=0.51, y radius=0.16];}
          \node[font=\tiny] at (35,-0.55) {Outdoor temperature ($^\circ$F)};
          \node[rotate=90, font=\tiny] at (0.2,6) {Gas used (therms)};
        \end{tikzpicture}
        \end{minipage}\hfill
        \begin{minipage}[b]{0.49\linewidth}\centering
        \begin{tikzpicture}[xscale=0.125, yscale=2.15]
          \draw[->, thick] (8,-0.8) -- (8,0.85);
          \draw[thick] (8,-0.8) -- (62,-0.8);
          \draw[dashed, slate] (8,0) -- (62,0);
          \foreach \y in {-0.6,-0.3,0,0.3,0.6} {\draw (7.5,\y) -- (8.5,\y);
            \node[left, font=\tiny] at (8,\y) {$\y$};}
          \foreach \x in {10,20,30,40,50,60} {\draw (\x,-0.825) -- (\x,-0.775);
            \node[below, font=\tiny] at (\x,-0.81) {\x};}
          \foreach \x/\e in {12/0.13, 15/-0.29, 18/0.38, 21/-0.14, 24/0.24, 27/-0.38, 30/0.00,
                             33/-0.52, 36/0.56, 39/-0.16, 42/0.32, 45/0.10, 48/-0.33, 51/0.15,
                             54/-0.07, 57/0.01}
            {\fill[navy] (\x,\e) circle[x radius=0.51, y radius=0.03];}
          \node[font=\tiny] at (35,-1.09) {Outdoor temperature ($^\circ$F)};
          \node[rotate=90, font=\tiny] at (-1.5,0) {Residual (therms)};
        \end{tikzpicture}
        \end{minipage}

        \vspace{4pt}
        {\centering\footnotesize\ttfamily
        \begin{tabular}{@{} l r r r r @{}}
        Predictor & Coef & SE Coef & T & P \\ \hline
        Constant & 11.69 & 0.2066 & 56.57 & 0.000 \\
        Temperature & -0.1597 & 0.00556 & -28.72 & 0.000 \\
        \end{tabular}\par
        S = 0.3075 \quad R-Sq = 98.3\% \quad R-Sq(adj) = 98.2\%\par}

        \vspace{6pt}
        \textbf{(a)} Describe the association between outdoor temperature and the natural gas
        used.\\[4pt]
        \gradedhere{\gradedblank[id=t27a, type=rubric, points=3, rubric=../rubrics/apstats-s01-practice-final-t27a.yaml]{\linewidth+6pt}{\gradedlinesheight{3}}}\ansline{Strong, negative, linear association: days with higher outdoor temperatures}\\
        \ansline{tend to have less natural gas used. There are no clear outliers.}\\
        \ansline{}\\
        \ansline{}\\
        \ansline{}\\
        \ansline{}\\
        \ansline{}\\
        \ansline{}

        \vspace{6pt}
        \boxguard[16]
        \textbf{(b)} Write the equation of the least-squares regression line, using the
        variable names. Then interpret its slope in context.\\[4pt]
        \gradedhere{\gradedblank[id=t27b, type=rubric, points=3, rubric=../rubrics/apstats-s01-practice-final-t27b.yaml]{\linewidth+6pt}{\gradedlinesheight{3}}}\ansline{predicted gas $= 11.69 - 0.1597$(temperature)}\\
        \ansline{For each additional degree Fahrenheit, the \emph{predicted} amount of natural gas}\\
        \ansline{used decreases by about 0.16 therms.}\\
        \ansline{}\\
        \ansline{}\\
        \ansline{}\\
        \ansline{}\\
        \ansline{}

        \vspace{6pt}
        \boxguard[21]
        \textbf{(c)} On one day the mean outdoor temperature was 36$^\circ$F and the house used
        6.5 therms of natural gas. Find the residual for this day and interpret it in context.
        \begin{work}
          \widehat{\text{gas}} &= 11.69 - 0.1597(36) = 5.94 \text{ therms} \\
          \text{residual} &= y - \hat{y} = 6.5 - 5.94 = 0.56 \text{ therms}
        \end{work}
        \strut\gradedhere{\gradedblank[id=t27c, type=rubric, points=3, rubric=../rubrics/apstats-s01-practice-final-t27c.yaml]{\linewidth+6pt}{\gradedlinesheight{3}}}\ansline{On this 36$^\circ$F day the house used about 0.56 therms \emph{more} natural gas than}\\
        \ansline{the line predicted.}\\
        \ansline{}\\
        \ansline{}\\
        \ansline{}\\
        \ansline{}\\
        \ansline{}\\
        \ansline{}

        \vspace{6pt}
        \boxguard[15]
        \textbf{(d)} Interpret the value of R-Sq in context.\\[4pt]
        \gradedhere{\gradedblank[id=t27d, type=rubric, points=2, rubric=../rubrics/apstats-s01-practice-final-t27d.yaml]{\linewidth+6pt}{\gradedlinesheight{3}}}\ansline{About 98.3\% of the variation in daily natural gas use is explained by the}\\
        \ansline{linear model with outdoor temperature.}\\
        \ansline{}\\
        \ansline{}\\
        \ansline{}\\
        \ansline{}\\
        \ansline{}\\
        \ansline{}

        \vspace{6pt}
        \boxguard[17]
        \textbf{(e)} Use the residual plot to decide whether a linear model is appropriate for
        these data. Then explain why the homeowner should not use this line to predict the
        natural gas used on a 90$^\circ$F summer day.\\[4pt]
        \gradedhere{\gradedblank[id=t27e, type=rubric, points=4, rubric=../rubrics/apstats-s01-practice-final-t27e.yaml]{\linewidth+6pt}{\gradedlinesheight{3}}}\ansline{Yes: the residuals are scattered randomly around 0 with no curved pattern, so a}\\
        \ansline{linear model is appropriate. But 90$^\circ$F is far above the 12--57$^\circ$F in the}\\
        \ansline{data. That is extrapolation, so the prediction is unreliable: the line gives}\\
        \ansline{about $-2.7$ therms, which is impossible.}\\
        \ansline{}\\
        \ansline{}\\
        \ansline{}\\
        \ansline{}

  % ── FRQ 28 ── Unit 3
  \boxguard[26]
  \item \textbf{(15 points)} A university dining service has 7{,}500 students on its meal
        plan. It wants to estimate the proportion of these students who leave campus on most
        weekends.

        \vspace{6pt}
        \textbf{(a)} Describe how to use a random number generator to choose a simple random
        sample of 250 of the students.\\[4pt]
        \gradedhere{\gradedblank[id=t28a, type=rubric, points=3, rubric=../rubrics/apstats-s01-practice-final-t28a.yaml]{\linewidth+6pt}{\gradedlinesheight{3}}}\ansline{Label the students 1 to 7{,}500. Use a random number generator to produce}\\
        \ansline{integers from 1 to 7{,}500, ignoring repeats, until 250 different labels are}\\
        \ansline{chosen. Survey the 250 students with those labels.}\\
        \ansline{}\\
        \ansline{}\\
        \ansline{}\\
        \ansline{}\\
        \ansline{}

        \vspace{6pt}
        \boxguard[17]
        \textbf{(b)} Staff knock on the dorm-room doors of the 250 selected students on
        Saturday afternoons, and reach only 100 of them. Name the bias this creates. Explain
        who is likely to be under-represented, and whether the estimate is likely too high or
        too low.\\[4pt]
        \gradedhere{\gradedblank[id=t28b, type=rubric, points=3, rubric=../rubrics/apstats-s01-practice-final-t28b.yaml]{\linewidth+6pt}{\gradedlinesheight{3}}}\ansline{Nonresponse bias. Students who leave campus on weekends are not in their rooms}\\
        \ansline{on Saturday afternoons, so they are under-represented. The estimated proportion}\\
        \ansline{who leave campus on most weekends is likely too low.}\\
        \ansline{}\\
        \ansline{}\\
        \ansline{}\\
        \ansline{}\\
        \ansline{}

        \vspace{6pt}
        \boxguard[18]
        \textbf{(c)} The dining service also wants to know whether calorie labels on the menu
        lower the calories students choose at lunch. One hundred twenty students volunteer; for
        one lunch, 60 will order from a menu with calorie labels and 60 from the same menu
        without labels, and the calories in each student's order will be recorded. Identify the
        experimental units, the treatments, and the response variable.\\[4pt]
        \gradedhere{\gradedblank[id=t28c, type=rubric, points=3, rubric=../rubrics/apstats-s01-practice-final-t28c.yaml]{\linewidth+6pt}{\gradedlinesheight{3}}}\ansline{Experimental units: the 120 student volunteers. Treatments: a menu with calorie}\\
        \ansline{labels, or the same menu without labels. Response variable: the calories in}\\
        \ansline{each student's lunch order.}\\
        \ansline{}\\
        \ansline{}\\
        \ansline{}\\
        \ansline{}\\
        \ansline{}

        \vspace{6pt}
        \boxguard[16]
        \textbf{(d)} Describe how to randomly assign the 120 volunteers to the two treatments,
        and explain why random assignment matters in this experiment.\\[4pt]
        \gradedhere{\gradedblank[id=t28d, type=rubric, points=3, rubric=../rubrics/apstats-s01-practice-final-t28d.yaml]{\linewidth+6pt}{\gradedlinesheight{3}}}\ansline{Write each volunteer's name on one of 120 identical slips, mix them in a hat, and}\\
        \ansline{draw 60: those students get the labeled menu, and the other 60 get the menu}\\
        \ansline{without labels. Random assignment balances other variables (such as appetite)}\\
        \ansline{between the groups, so a difference in calories can be attributed to the labels.}\\
        \ansline{}\\
        \ansline{}\\
        \ansline{}\\
        \ansline{}

        \vspace{6pt}
        \boxguard[17]
        \textbf{(e)} The students who saw the calorie labels chose significantly fewer calories,
        on average, than the students who did not. What can the dining service conclude about
        cause and effect, and to which students does the conclusion apply? Explain.\\[4pt]
        \gradedhere{\gradedblank[id=t28e, type=rubric, points=3, rubric=../rubrics/apstats-s01-practice-final-t28e.yaml]{\linewidth+6pt}{\gradedlinesheight{3}}}\ansline{The treatments were randomly assigned and the difference is significant, so the}\\
        \ansline{dining service can conclude the labels caused students to choose fewer calories.}\\
        \ansline{But the subjects were volunteers, not a random sample, so the conclusion applies}\\
        \ansline{only to students like these volunteers, not to all 7{,}500 on the meal plan.}\\
        \ansline{}\\
        \ansline{}\\
        \ansline{}\\
        \ansline{}

\end{enumerate}

\end{document}
